% -*-mode: LaTeX; coding: utf-8; -*-
\documentclass{article}
\usepackage[a4paper, margin=2cm]{geometry}
% For vakker crypto.  Dokumentasjonen ligger i katalogen
\usepackage{aspen}
% For vakrere referanser
\usepackage{cite}


\title{Zero-knowledge proof}
\author{[TaSK]}
\date{juni 2026}
\begin{document}
\maketitle

   
\section{Innledning}
Dette tekniske notatet gir en oversikt over en publisert
artikkel~\cite{Goldwasser1985Knowledge,Goldwasser1989Knowledge}.
Fordi dette notatet beskriver en artikkel er her ingen nye resultater.
Termer som ``artikkelen'', ``kapittel 2'', ``figur 2'', og tilsvarende
refererer til~\cite{Goldwasser1989Knowledge}.  Der det er nødvendig å
referere til denne teksten bruker vi eksplisitt ``i dette notatet''.

Det gjøres ingen antagelser om kanaler som meldinger sendes på, utover
ordning i tid (en melding ankommer ikke føær den er sendt).  I
særdeleshet antas det ikke at kanaler tilbyr konfidensialitet. 

%%%%%%%%%%%%%%%%%%%%%%%%%%%%%%%%%%%%%%%%%%%%%%%%%%%%%%%%%%%%%%%%%%%%%%%%%%%%%%%
\section{Bakgrunn}
Tema er hvordan  Alice kan overbevise Bob at hun kjenner en
hemmelighet, uten å hverken fortelle noe av verdi til Bob (ut over at
hun kjenner hemmeligheten), og heller ikke til Charlie (som ikke
kjenner hemmeligheten) som lytter på samtalen.

For eksempel: Gitt at Alice har konfidensielle data, for eksempel
et dokument; la oss kalle dokumentet for $D$.  Hun tilbyr ett eller
annet til andre som også har dokumentet.  Bob mener han har det, og
ønkser å bevise til Alice om det.  Når følgende protokoll er ferdig er
\APri{A} overbevist om at \APri{B} også har (tilgang til) $D$:

\begin{ASteps*}
  \ASend*{A}{B}{\ANonce{A}} \\
  \ASend*{B}{A}{\AChash{\ANonce{A} \AConcat D}}
\end{ASteps*}
%
hvor notasjonen $\AConcat$ betyr sammensetting
(\textit{concatination}), $\ANonce{A}$ ``et svært stort tilfeldig
tall'' (la oss si $\approx2^{256}$), og $\AChash{m}$ et digitalit
fingeravtrykk (\textit{hashing}) av $m$.

Fordi han lytter på all kommunikasjon kjenner Charlie nå $\ANonce{A}$
og $\AChash{\ANonce{A} \AConcat D}$.  Det vil si at om han har stjålet et
dokument $D^\prime$ kan han undersøke om $D^\prime = D$ ved å beregne
$\AChash{\ANonce{A} \AConcat D^\prime}$ og sammenligne med
$\AChash{\ANonce{A} \AConcat D}$.  Det vil at protokollen mellom
$\APri{A}$ og $\APri{B}$ lakk (meta)informasjon om $D$ til
$\APri{C}$.

Evnen til å kunne verifisere om $D=D^\prime$ er en betydelig forde som
kan gi næring til mange typer angrep.  \textit{Zero Knowledge} handler
om å overbevise den andre uten å lekke informasjon. 


\section{Matematisk utgangspunkt}
Den matematiske konteksten i artikkelen er språk og
kompleksitetsklasser; vi gir derfor en kort oppfriskning.

\subsection{Språk}
Alan Turing var den første til å representere (matematiske) spørsmål
med symboler (tekst)~\cite{Turing1937Computable} men det var først 25 år
senere at det ble forstått at dersom et problem representeres som en
kompakt streng (sekvens av symboler) så vil lengden av strengen være
et mål på ``størrelsen'' av problemet~\cite{Hartmanis1965Complexity}.
For en introduksjon til anbefaler vi~\cite{AroraBarak2009Complexity};
første kapittel går nøye gjennom representasjon av problemer.

I spørsmål om involverer tall er det åpenbart hvordan problemet skal
representeres. For eksempel: Om spørsmålet er: ``Er tallet et
primtall?'' så er problemet som skal besvares et tall.
Tallet (problemet) kan være $X=323$ (som vil si $(101000011)_2$).
Lengden på dette problemet er $9$.

Når spørsmålet omhandler en graf, for eksempel om grafen er forbundet,
kan kan vi konvertere grafen til en matrise (hvor $M_{ij}=0$ om det
ikke er en kant mellom nodene $i$ og $j$).  Når en slik matrise
skrives ut blir resultatet igjen en streng (av $0$ og $1$).

Når vi bruker strenger til å representere problemer på denne måten kan
vi betrakte strengene som et ``ord''; hver streng representerer ett
(og bare ett) problem.  Indirekte har vi nå også fått et mål på
størrelse: Jo lenger streng jo større problem.

Om vi ser på de to uendelige settene av alle probløemer (tall eller
fgrafer) så vil spørsmålet (``Er tallet prim?'', eller ``Er grafen
forbundet?'') dele  settet i to: Ja eller nei; de begge tilfellene vil
to sub-settene ikke være like store.  Settet som representerer ``ja'',
det er ord i \emph{et språk}.

Spørsmålet om $X$ er prim, det heter \textsc{Primes}.  Den uformelle
beskrivelksen blir at gitt en binær streng som vi tolker som et
tall $X$ slik at $X\in\textbf{N}$, er $X$ et primtall?  Her behøver vi
heldigvis ikke definere hva primtall er, for det vet vi.  Språket
\textsc{Primes} er settet med (binære) strenger som representerer tall
som er prim.
Akkurat som ikke alle bokstavkombinasjoner tilhører språket norsk,
tilhører ikke alle tall språket \textsc{Primes}.

Spørsmålet om en graf inneholder en komplett sub-graf (klikk) det
heter \textsc{Clique}.  Språket $L_\textsc{Clique}$ er ``ord'' hvor
svaret er ``ja'':
\begin{equation}
  L_{\textsc{CLIQUE}} = \{ \langle G, k \rangle \mid G \text{ har en
    klikk med størrelse } k \}
\end{equation}
Her er $G$ en streng som representerer en graf og $k$ et tall
(størrelsen på klikken).



\subsection{Kompleksitesklasser}

\subsubsection{Klassen P}
De språkene som kan konstrueres i tid og rom som ikke vokser mer enn
polynomisk relativ til lengden på spørsmålet, de sies å tilhøre
klassen \textbf{P}.

For eksempel \textsc{Partall}:
\begin{equation}
  \textsl{PARTALL}: \{n: n \textrm{er partall}
\end{equation}
Algoritmen for å bestemme om et tall er i \textsc{Partall} er å se om
siste bit er 0 eller 1.  Det tar fast tid (uavhengig av lengen) og
derfor klart mindre enn polynomisk tid.

Gitt en graf $G$ kan vi bruke bredde-først eller dybde-først for å
finne om én node er forbundet med en annen.  Ingen av dem vokser mer
enn polynomisk når størrelsen på $G$ er målet.  Det vil si at språkket
\textsc{Sti} er i \textbf{P}:
\begin{equation}
  \textrm{STI} = \{(G, a, b): \textrm{Det er en sti fra $a$ til $b$ i
    $G$} 
\end{equation}

%
%
\subsubsection{Klassen NP}

En annen klasse, som inkluderer settet \textbf{P}, er \textbf{NP}.
Språk i denne klassen er de hvor det er mulig å verifisere et
bevis\footnote{På engelsk har \emph{proof} og \emph{evidence}
forskjellig betydning.  Når vi snakker om \emph{bevis} her mener vi
\emph{evidence}.} på at en instans har ``ja''-svar i polynomisk tid.

For å se at $P$ er innholdt i $NP$ er det tilstrekkelig å notere at
elementene  i $P$ kan konstrueres i polynomisk tid.  Det vil si det
ikke er behov for noe bevis i det hele tatt; man kan bare regne ut
løsningen.

Et språk som ofte brukes som eksempel er \textsc{Sammensatt}:
\begin{equation}
  \textrm{SAMMENSATT}: \{n: \textrm{n er sammensatt (ikke prim)}\}
\end{equation}
Et bevis for at $n$ er sammendatt (et ``ja''-svar) er $n=xy$.
Kompleksiteten ved å verifisere, altså multiplisere, øker med lengden
på $x$ og $y$ (som naturligvis øker med lengden på $n$), men kun
polynomisk.  Altså er \textsc{Sammensatt} i \textbf{NP}.  Kjernen er
at det som skal gjøres er ikke å finne ut om et tall er sammensatt,
men undersøke om et bevis er gyldig.

\textsc{Primes} er også i \textbf{NP} fordi det finnes et bevis somk
kan verifiseres i polynomisk tid.  Dersom $n$ er prim da vil det
finnes en $a^{n-1} = 1 (= n)$ og for hver primfaktor $q$ av
$n-1$ vil $a^{(n-1)/q} \not= 1 (= n)$.  Et bevis for at $n$ er
prim vil være $a$ sammen med faktoriseringen av $n-1$ (og tilsvarende
bevis for at de faktorene er prim).

Her er det viktig å huske at å finne ut om $n$ er prim er ett problem
mens avgjøre om beviset på at $n$ er prim er gyldig, er et annet.
Kravet for at et problem er i \textbf{NP} er altså at det finnes et
bevis (og ikke nødvendigvis at beviset er ``lett'' å finne).



%%%%%%%%%%%%%%%%%%%%
\section{Notasjon}

\bibliographystyle{plain}
\bibliography{zero}

\end{document}





